\documentclass[aspectratio=169,11pt]{beamer}

\usetheme{Madrid}
\usecolortheme{default}
\usefonttheme{professionalfonts}
\setbeamertemplate{navigation symbols}{}
\setbeamertemplate{footline}[frame number]

\usepackage{amsmath, amssymb, booktabs, graphicx}

\newcommand{\E}{\mathbb{E}}
\newcommand{\1}{\mathbf{1}}
\newcommand{\Var}{\operatorname{Var}}
\newcommand{\argmin}{\operatorname*{arg\,min}}
\newcommand{\supp}{\operatorname{supp}}

\title{High-Dimensional Controls, Heterogeneity, and Double/Debiased ML}
\subtitle{Empirical Methods --- Week 10}
\author{Teaching deck draft}
\date{}

\begin{document}

\begin{frame}
  \titlepage
\end{frame}

\begin{frame}{Week 10 question}
\begin{itemize}
    \item How can flexible methods help estimate causal effects without turning nuisance prediction into the answer?
    \item When are high-dimensional controls a design asset rather than a post-selection problem?
    \item How do orthogonal scores and cross-fitting protect inference?
    \item When is treatment-effect heterogeneity substantively meaningful?
\end{itemize}

\medskip
Prediction accuracy is useful only after the estimand, timing, overlap, and identifying assumptions are explicit.
\end{frame}

\begin{frame}{Why high-dimensional controls create new problems}
\small
\begin{itemize}
    \item Applied data contain worker histories, firms, places, text measures, lagged outcomes, and interactions.
    \item Including everything can overfit, create collinearity, and weaken precision.
    \item Selecting controls after seeing results invalidates ordinary post-selection inference.
    \item Hand-picked controls can omit important confounders that predict treatment.
\end{itemize}

\medskip
\[
(Y_i(1),Y_i(0)) \perp D_i \mid X_i,
\qquad
0<\Pr(D_i=1\mid X_i)<1
\]

\medskip
The algorithm helps use rich observables. It does not make the observables sufficient.
\end{frame}

\begin{frame}{Post-double selection}
\small
\[
Y_i=\theta_0D_i+g_0(X_i)+U_i,
\qquad
D_i=m_0(X_i)+V_i
\]

\[
\widehat S_Y=\supp(\widehat\gamma_Y),
\qquad
\widehat S_D=\supp(\widehat\gamma_D),
\qquad
\widehat S=\widehat S_Y\cup \widehat S_D
\]

\begin{itemize}
    \item Use lasso-style selection in the outcome equation and treatment equation.
    \item Estimate the treatment effect using the union of selected controls.
    \item Best when approximate sparsity is plausible and the target is low-dimensional.
    \item Does not solve unobserved confounding, nonlinear misspecification, or poor overlap.
\end{itemize}
\end{frame}

\begin{frame}{Orthogonalization: DML math}
\small
\[
g_0(X_i)=\E[Y_i\mid X_i],
\qquad
m_0(X_i)=\E[D_i\mid X_i]
\]

\[
\widetilde Y_i=Y_i-g_0(X_i),
\qquad
\widetilde D_i=D_i-m_0(X_i)
\]

\[
\psi(W_i;\theta,\eta)=
\left(D_i-m(X_i)\right)
\left[Y_i-g(X_i)-\theta\left(D_i-m(X_i)\right)\right],
\qquad
\eta=(g,m)
\]

\[
\widehat\theta_{DML}
=
\left(\frac{1}{n}\sum_i\widehat{\widetilde D}_i^2\right)^{-1}
\left(\frac{1}{n}\sum_i\widehat{\widetilde D}_i\widehat{\widetilde Y}_i\right)
\]

Orthogonality makes small nuisance errors second-order for the target score.
\end{frame}

\begin{frame}{Cross-fitting}
\small
Split the sample into folds \(I_1,\ldots,I_K\). For each held-out fold:

\[
\widehat{\widetilde Y}_i
=
Y_i-\widehat g^{(-k)}(X_i),
\qquad
\widehat{\widetilde D}_i
=
D_i-\widehat m^{(-k)}(X_i),
\qquad
i\in I_k
\]

\medskip
\begin{itemize}
    \item Train nuisance learners on \(I_k^c\), evaluate scores on \(I_k\).
    \item Rotate across folds and stack held-out residuals.
    \item Tune hyperparameters inside training folds.
    \item Cluster folds when the identifying variation or errors are clustered.
\end{itemize}

\medskip
Cross-fitting protects the target step from overfit nuisance predictions.
\end{frame}

\begin{frame}{What DML identifies}
\small
\begin{tabular}{p{0.24\linewidth} p{0.32\linewidth} p{0.32\linewidth}}
\toprule
Object & Helps when & Does not fix \\
\midrule
High-dimensional controls & many pre-treatment confounders are observed & omitted unobservables or post-treatment controls \\
Orthogonal score & nuisance functions are complex but auxiliary & a wrong estimand or weak design \\
Cross-fitting & overfit nuisance estimates threaten inference & tuning leakage or bad support \\
DML estimate & target is low-dimensional and design is credible & extrapolation outside overlap \\
\bottomrule
\end{tabular}

\medskip
Report the estimand, identifying assumption, nuisance learners, fold plan, overlap checks, and learner sensitivity.
\end{frame}

\begin{frame}{CATEs and causal forests}
\small
\[
\tau(x)=\E[Y_i(1)-Y_i(0)\mid X_i=x]
\]

\[
\widehat\tau(x)=
\argmin_{\tau}
\sum_i \alpha_i(x)
\left[
\left(Y_i-\widehat\mu(X_i)\right)
-
\tau\left(D_i-\widehat e(X_i)\right)
\right]^2
\]

\begin{itemize}
    \item Forest weights \(\alpha_i(x)\) define local comparison groups.
    \item Honesty separates split selection from effect estimation.
    \item Useful for heterogeneity discovery, targeting diagnostics, and mechanism hypotheses.
    \item Variable importance is not a causal mechanism.
\end{itemize}
\end{frame}

\begin{frame}{Implementation pitfalls}
\small
\begin{tabular}{p{0.25\linewidth} p{0.37\linewidth} p{0.26\linewidth}}
\toprule
Pitfall & Why it matters & Check \\
\midrule
Weak overlap & flexible learners extrapolate in thin support & propensity distribution, trimming \\
Feature leakage & post-treatment features encode the answer & timing diagram \\
Learner choice & nuisance bias/variance affects target stability & lasso, ridge, forest sensitivity \\
Tuning leakage & full-sample tuning weakens splitting discipline & nested tuning \\
Fold instability & estimates depend on random split & seed and fold diagnostics \\
Noisy heterogeneity & CATEs overfit sparse subgroups & honest forests, hold-out validation \\
\bottomrule
\end{tabular}
\end{frame}

\begin{frame}{Research Lab design}
\begin{columns}[T,onlytextwidth]
\begin{column}{0.32\linewidth}
\textbf{Reproduce}
\begin{itemize}
    \item synthetic online labor data
    \item baseline regression
    \item post-double selection
    \item cross-fitted DML
\end{itemize}
\end{column}
\begin{column}{0.32\linewidth}
\textbf{Diagnose}
\begin{itemize}
    \item overlap support
    \item nuisance \(R^2\)
    \item leakage audit
    \item fold stability
\end{itemize}
\end{column}
\begin{column}{0.32\linewidth}
\textbf{Transfer}
\begin{itemize}
    \item youth program data
    \item CATE by baseline risk
    \item honest-tree heterogeneity
    \item uncertainty memo
\end{itemize}
\end{column}
\end{columns}

\medskip
The lab reproduces design logic, not official paper estimates.
\end{frame}

\begin{frame}{Takeaways}
\begin{itemize}
    \item Start with the economic question and target parameter.
    \item High-dimensional controls help only when the relevant confounders are observed and timed correctly.
    \item Orthogonalization protects inference from nuisance estimation error, not from bad identification.
    \item Cross-fitting is part of the estimator, not a cosmetic validation step.
    \item Causal forests organize heterogeneity, but support, uncertainty, and economics determine interpretation.
    \item Report these tools as research-design infrastructure, not as black-box proof.
\end{itemize}
\end{frame}

\end{document}
